\documentclass[11pt]{article}
\usepackage{amsmath,amssymb,amsthm}
\usepackage{geometry}
\usepackage{hyperref}
\geometry{margin=2.6cm}
\newtheorem{theorem}{Theorem}
\title{Local bifurcation at the first Dirichlet eigenvalue}
\author{}
\date{}
\begin{document}
\maketitle

Let \(\Omega\subset\mathbb R^n\) be bounded and connected with
\(C^{2,\alpha}\) boundary, \(0<\alpha<1\). Let \(\lambda_1\) be the first
Dirichlet eigenvalue of \(-\Delta\), and let \(\phi_1>0\) be a corresponding
eigenfunction. Write
\[
 N=\int_\Omega\phi_1^2\,dx>0,\qquad
 A_4=\int_\Omega\phi_1^4\,dx>0.
\]
The formula below allows any normalization of \(\phi_1\).

\begin{theorem}
For the problem
\[
 -\Delta u=\lambda u+u^3\quad\text{in }\Omega,\qquad
 u|_{\partial\Omega}=0,
\]
there are \(\varepsilon>0\) and smooth maps
\(s\mapsto\lambda(s)\) and \(s\mapsto u(s)\in X\), where
\[
 X=\{u\in C^{2,\alpha}(\overline\Omega):u|_{\partial\Omega}=0\},
\]
such that \((\lambda(0),u(0))=(\lambda_1,0)\), \(u(s)\ne0\) for
\(0<|s|<\varepsilon\), and
\[
 u(s)=s\phi_1+O(s^3)\quad\text{in }X,\qquad
 \lambda(s)=\lambda_1-\frac{A_4}{N}s^2+O(s^4).
\]
Thus \((\lambda_1,0)\) is a local bifurcation point, and
\(\lambda(s)<\lambda_1\) for sufficiently small \(s\ne0\).
Moreover \(\lambda(-s)=\lambda(s)\), \(u(-s)=-u(s)\), and every
nontrivial solution in a sufficiently small neighborhood of
\((\lambda_1,0)\) in \(\mathbb R\times X\) belongs to this curve.
\end{theorem}

\begin{proof}
Put \(Y=C^{0,\alpha}(\overline\Omega)\) and
\[
 F(\lambda,u)=-\Delta u-\lambda u-u^3:
       \mathbb R\times X\longrightarrow Y.
\]
This polynomial map is \(C^\infty\), and \(F(\lambda,0)=0\).
Schauder theory makes the Dirichlet operator
\(A=-\Delta:X\to Y\) an isomorphism. The inclusion
\(J:X\hookrightarrow Y\) is compact; hence
\(L=D_uF(\lambda_1,0)=A-\lambda_1J\) is Fredholm of index zero.
The principal Dirichlet eigenvalue theorem on a connected domain gives
\(\ker L=\operatorname{span}\{\phi_1\}\).
Integration by parts gives
\(\int_\Omega\phi_1 Lv\,dx=0\) for every \(v\in X\).
Because \(\operatorname{Range}L\) has codimension one,
\[
 \operatorname{Range}L
 =\left\{f\in Y:\int_\Omega\phi_1 f\,dx=0\right\}. \tag{1}
\]
Also \(D_{\lambda u}^2F(\lambda_1,0)[\phi_1]=-\phi_1\) is outside
this range, since its pairing with \(\phi_1\) is \(-N\).

Set
\[
 W=\left\{v\in X:\int_\Omega\phi_1v\,dx=0\right\},\qquad
 Qf=f-\frac{\int_\Omega\phi_1f\,dx}{N}\phi_1.
\]
Then \(X=\operatorname{span}\{\phi_1\}\oplus W\),
\(QY=\operatorname{Range}L\), and \(L|_W:W\to QY\) is an
isomorphism. Write \(u=s\phi_1+w\), where \(w\in W\).
The Banach-space implicit function theorem applied to
\[
 G(s,\lambda,w)=QF(\lambda,s\phi_1+w)=0
\]
at \((0,\lambda_1,0)\) gives a unique smooth
\(w=w(s,\lambda)\in W\) near that point.

Since \(G(0,\lambda,0)=0\), uniqueness gives
\(w(0,\lambda)=0\) for every nearby \(\lambda\). Moreover,
\[
 G_s(0,\lambda,0)
 =Q[(\lambda_1-\lambda)\phi_1]=0.
\]
The operator \(G_w(0,\lambda,0)\) remains invertible for nearby
\(\lambda\), so implicit differentiation gives
\(w_s(0,\lambda)=0\). The oddness
\(F(\lambda,-u)=-F(\lambda,u)\) and uniqueness give
\(w(-s,\lambda)=-w(s,\lambda)\). Thus \(w\) is odd in \(s\),
and its zeroth, first and second derivatives at zero vanish.
Uniform Taylor expansion yields
\[
 \|w(s,\lambda)\|_X\le C|s|^3                       \tag{2}
\]
for \(\lambda\) in a sufficiently small interval about \(\lambda_1\).

Define the unnormalized scalar equation
\[
 g(s,\lambda)
 =\int_\Omega\phi_1F(\lambda,s\phi_1+w(s,\lambda))\,dx.
\]
By (1), \(w\in W\), and integration by parts,
\(\int_\Omega\phi_1(-\Delta-\lambda)w\,dx=0\).
Expanding the cubic term and using (2) gives, uniformly near
\(\lambda_1\),
\[
 g(s,\lambda)
 =(\lambda_1-\lambda)Ns-A_4s^3+O(s^5).          \tag{3}
\]
In particular the cross term \(s^2w\) is \(O(s^5)\).

The smooth function \(g\) is odd in \(s\). Define
\[
 H(s,\lambda)=\int_0^1\partial_sg(ts,\lambda)\,dt.
\]
Then \(g(s,\lambda)=sH(s,\lambda)\), and \(H\) is smooth and even
in \(s\). Equation (3) gives
\[
 H(s,\lambda)=(\lambda_1-\lambda)N-A_4s^2+O(s^4),
 \qquad
 \partial_\lambda H(0,\lambda_1)=-N\ne0.
\]
The scalar implicit function theorem therefore gives a unique smooth
\(\lambda=\lambda(s)\) near zero such that
\(H(s,\lambda(s))=0\) and \(\lambda(0)=\lambda_1\).
Evenness of \(H\) and uniqueness give
\(\lambda(-s)=\lambda(s)\). Substitution yields
\[
 \lambda(s)=\lambda_1-\frac{A_4}{N}s^2+O(s^4).
\]
Since \(A_4/N>0\), the branch lies in \(\lambda<\lambda_1\)
for small \(s\ne0\). Set
\(u(s)=s\phi_1+w(s,\lambda(s))\).
By (2), \(u(s)=s\phi_1+O(s^3)\), so \(u(s)\ne0\) when
\(s\ne0\) is sufficiently small.
Oddness of \(w\) gives \(u(-s)=-u(s)\).

Finally, any nearby solution decomposes uniquely as
\(u=s\phi_1+w\), \(w\in W\). The auxiliary equation forces
\(w=w(s,\lambda)\). If \(s=0\), then \(w=w(0,\lambda)=0\).
For \(s\ne0\), the remaining scalar equation \(g(s,\lambda)=0\)
is equivalent to \(H(s,\lambda)=0\), whose local uniqueness
forces \(\lambda=\lambda(s)\). Hence every nearby nontrivial
solution lies on the stated curve.
\end{proof}

\begin{thebibliography}{9}
\bibitem{CR71}
M.~G. Crandall and P.~H. Rabinowitz,
\emph{Bifurcation from simple eigenvalues},
Journal of Functional Analysis \textbf{8} (1971), 321--340.
\href{https://doi.org/10.1016/0022-1236(71)90015-2}{doi:10.1016/0022-1236(71)90015-2}.
\end{thebibliography}
\end{document}
